\chapter{Материалы, реализация и границы проверки}
\label{chap:materials}

\section{Как связаны текст, код и результаты}
Все результаты в этом описании относятся к единому вычислительному
конвейеру. Модуль \file{src/physics/} задаёт гамильтониан, каналы и
эволюцию матрицы плотности; \file{src/blocks/observe.py} строит
наблюдаемые; \file{src/blocks/target.py} определяет первое пересечение
порога; \file{src/data/windows.py} разрезает траектории на окна; модули
\file{src/models/} реализуют сравниваемые прогнозы; \file{src/evaluation/}
вычисляет метрики. Скрипты \file{experiments/} соединяют эти части в
именованные E- и R-опыты, а их строки результатов сохраняются вместе с
параметрами запуска в \file{experiments/results/}.

Такое разделение позволяет читать текст от физической постановки к
вычислимому результату. Сначала выбираются $H$, $L_k$, параметры начального
состояния и профиль среды; затем решается уравнение ГКСЛ, строится ряд
наблюдаемых и определяется время события; только после этого выбираются
окна, разбиение на train/validation/test и модель. Каждая строка каталога
экспериментов в приложении указывает, какая именно часть этой цепочки
меняется. Поэтому, например, R15 проверяет устойчивость к шуму входа, R21
меняет саму скорость среды, R23 меняет доступную информацию о двух кубитах,
а R28 меняет физический канал диссипации.

Основные статьи используют точные ансамблевые средние, вычисленные из
$\rho(t)$. Поздние R29 и R30 моделируют конечное число повторных
приготовлений и строят производные признаки из тех же исходов. Это разные
уровни постановки: результаты с точными средними служат контролем алгоритма,
а конечные измерения проверяют чувствительность конвейера к статистике
наблюдения. Во всех этих опытах метка $T_2$ вычислена по идеальной
моделируемой траектории.

\section{Заметка: нормировка входа BiLSTM}
\label{sec:bug}
В ранней версии кода окна при обучении BiLSTM (гибридной модели и BiLSTM
без оценки) подавались в сеть без нормировки, а при предсказании
нормировались (раздел~\ref{chap:learning}). Из-за этого BiLSTM без оценки на
одном кубите давал $R^2$ на уровне константы. Ошибка исправлена, все
результаты рекуррентных моделей в этом описании получены после исправления;
одинаковую нормировку на обоих этапах проверяет автоматический тест.

\section{Использованные материалы}
Теоретическая часть опирается на уравнение ГКСЛ и определения когерентности,
чистоты и concurrence, приведённые в первых главах. Для постановки измерений
и интерпретации результатов использованы работы о прогнозе по повторным
приготовлениям и о последовательностях квантовых измерений, приведённые в
обзоре литературы. Микроскопическая глава R28 воспроизводит формулы и
опубликованные рисунки работ Delgado с соавторами; оцифрованные точки служат
проверкой формы кривых, а не самостоятельным источником данных.

Числа в таблицах взяты из сохранённых CSV. Если один и тот же вопрос
проверялся несколькими seed, в тексте приводится среднее и разброс; если
значение относится к одному E-запуску, это прямо указано. Два файла E1
относятся к разным запускам и описываются раздельно. Архивные E9, E10,
E12--E14 удалены из рабочего набора и не используются для выводов.

\section{Следующие проверки}
\begin{enumerate}
\item Пересчитать E5 после исправления привязки окна к своей траектории;
      прежняя таблица обратного метода сохраняется как исторический результат,
      но не используется для окончательного сравнения.
\item Повторить R30 для локального входа после исправления BiLSTM. Затем
      сравнить шесть локальных средних и томографию при одном фиксированном
      протоколе разбиения и обучения.
\item Построить бутстреп R13 по явному идентификатору траектории, а не по
      совпадающему численному значению $T_2$.
\item Расширить R29/R30 на коррелированный шум и зафиксировать аппаратные
      длительности приготовления, эволюции и считывания, чтобы перевести
      бюджет числа исходов во время.
\item Для последовательных измерений одного экземпляра обучить генератор
      по правдоподобию записи, сравнить его с прямым прогнозом и проверить
      идентифицируемость дрейфа $\gamma(t)$ при выбранных осях измерения.
\item В Kondo-модели включить сдвиг уровней в динамику и сопоставить
      секулярный генератор с несекулярным расчётом в области близких
      боровских частот.
\end{enumerate}
